\documentclass[11pt]{article}

% ---------- packages ----------
\usepackage[utf8]{inputenc}
\usepackage[T1]{fontenc}
\usepackage{amsmath, amssymb, amsthm}
\usepackage[dvipsnames, svgnames]{xcolor}
\usepackage{graphicx}
\usepackage{hyperref}
\usepackage{booktabs}
\usepackage{microtype}
\usepackage{geometry}
\usepackage{enumitem}
\geometry{a4paper, margin=2.5cm}
\usepackage{microtype}
\usepackage{amsmath,amssymb,amsthm}
\usepackage{mathtools}
\usepackage{bm}
\usepackage{graphicx}
\usepackage{booktabs}
\usepackage{hyperref}
\usepackage{xcolor}
\usepackage{algorithm}
\usepackage{algpseudocode}
\usepackage{geometry}
\geometry{margin=2.5cm}
\usepackage{subcaption}
\usepackage{float}

\usepackage{fontawesome5}
\usepackage{hyperref}

\makeatletter
\newcommand{\github}[1]{%
   \href{#1}{\faGithubSquare}%
}
\makeatother

% ---------- hyperref config ----------
\hypersetup{
  colorlinks = true,
  linkcolor  = blue!60!black,
  citecolor  = teal,
  urlcolor   = teal,
  pdftitle   = {Entropic Semantic Diffusion: A Spectral Latent Diffusion Model},
  pdfauthor  = {Lorenzo Moriondo},
}

% ---------- theorem environments ----------
\newtheorem{definition}{Definition}
\newtheorem{proposition}{Proposition}
\newtheorem{remark}{Remark}
\newtheorem{conjecture}{Conjecture}

% ---------- custom commands ----------
\newcommand{\ESDM}{\textsc{esdm}}
\newcommand{\ALDSC}{\textsc{ald-sc}}
\newcommand{\VDT}{\textsc{vdt}}
\newcommand{\Lf}{L_F}
\newcommand{\Uq}{U_q}
\newcommand{\Lq}{\Lambda_q}
\newcommand{\ed}{\lambda^{\mathrm{ED}}}
\newcommand{\R}{\mathbb{R}}

% ============================================================
\begin{document}
% ============================================================

\title{%
  \textbf{Entropic Semantic Diffusion:\\ \ALDSC --- A Spectral Latent Diffusion Model}
}

\author{%
  Lorenzo Moriondo \\[2pt]   \small tuned.org.uk\\   \small ORCID: 0000-0002-8804-2963\\   \texttt{tunedconsulting@gmail.com} }  \date{July 2026 \\ 
    \href{https://github.com/tuned-org-uk/entropic-semantic-diffusion}{%
     \textsc{\Large\faGithub} tuned-org-uk/entropic-semantic-diffusion} \\
     \href{https://github.com/tuned-org-uk/arrowspace-latent-diffusion}{%
     \textsc{\Large\faGithub} tuned-org-uk/arrowspace-latent-diffusion}
  }
    \maketitle  % ============================================================
    \begin{abstract} %============================================================  
    Most diffusion models generate images by denoising a latent or pixel-space state according to an externally prescribed noise schedule. While this design is effective, it leaves the semantic structure of the latent space largely unconstrained: smoothness, coherence, and global feature organisation emerge only indirectly from reconstruction and denoising losses. In this work we reframe diffusion generation through the ArrowSpace dual-space hypothesis, according to which every representation space is simultaneously an item-space and a feature-space endowed with its own graph geometry. We propose the \emph{Entropic Semantic Diffusion Model} (\ESDM), a spectrally conditioned latent diffusion architecture that combines a conventional spatial latent with a frozen feature-space geometric prior built from an ArrowSpace graph Laplacian.  The model is organised around a 2.5-D latent representation. A VAE-style encoder maps an image to (i) a spatial latent tensor carrying local visual detail and (ii) a feature field whose channels are projected onto the smooth eigenspace of a corpus-level ArrowSpace Laplacian. From this projected feature field, the model extracts a compact spectral chart consisting of low-mode coordinates, band-energy allocations, and an energy-dispersion prior. A latent diffusion backbone then denoises only the spatial latent, while a topology-adaptive decoder reconstructs the image under spectral conditioning from the frozen feature geometry. In this way, local image synthesis remains computationally efficient but global semantic organisation is explicitly constrained by a feature-space manifold.  Based on the work made with \VDT~and \ESDM, the \emph{Arrowspace Latent Diffusion with Spectral Chart Conditioning} (\ALDSC)  model is developed as the minimal viable implementation of this approach. 
    \end{abstract}  
    
    \vspace{1em} \noindent \textbf{Keywords:} diffusion models, ArrowSpace, graph Laplacian, spectral conditioning, latent diffusion, semantic search, energy dispersion, thermal time, spectral decoding.  \newpage \tableofcontents \newpage  % ============================================================
    \section{Introduction} \label{sec:introduction} % ============================================================  
    Latent diffusion models achieve high-quality image synthesis by moving the iterative denoising process from pixel space to a compressed latent space. Refactoring part ofthe work introduced in \cite{Moriondo2026VDT}, instead of denoising on the image manifold itself, the model denoises on an autoencoder latent whose global structure is usually only weakly constrained. The practical question is therefore not only how to compress images but how to ensure that the compressed space preserves coherent semantic organisation.  This paper approaches that question through the ArrowSpace dual-space view \cite{Moriondo2025JOSS, MoriondoAzizi2026}. In ArrowSpace, a vector space is not merely a collection of items embedded in a metric ambient space; it is also a feature-space whose internal relations can be described by a graph Laplacian and corresponding spectral structure. The resulting geometry does not replace ordinary item-space similarity but conditions it: item vectors may be projected, filtered, or interpreted relative to a feature-space manifold. Inspired by Barontini's experimental construction of entropic time in a closed cold-atom setting \cite{Barontini2026}, \ESDM~design an entropic clock in which time, diffusion schedule, and semantic diffusion should all emerge following the internal thermodynamics of a two-sector system. The latest version of the model demonstrate the use of the frozen feature-space prior to define its valid semantic subspace; its decoder reconstruct on the projection of the item-space into the feature-space rather than on an unconstrained ambient sphere.  Accordingly, the present paper reformulates \ESDM~as a topology-adaptive latent diffusion model with three commitments: 
    \begin{enumerate}[leftmargin=1.5em] 
      \item a frozen ArrowSpace prior $(\Lf,\Uq,\Lq,\ed)$ built once from corpus feature statistics;
      \item a "2.5-D" (informed by discrete topology) latent consisting of a standard spatial latent plus a compact spectral chart;
      \item a decoder conditioned by feature-space spectral geometry rather than only by Euclidean latent coordinates. 
    \end{enumerate}  
    This yields a model that is recognisably similar to modern latent diffusion systems but whose semantic organisation is explicitly governed by graph Laplacian structure.  
    % ============================================================
    
    \section{Background}
    \label{sec:background}
    
    \subsection{ArrowSpace and feature-space geometry}
    \label{sec:arrowspace} 
    ArrowSpace introduces a graph-based treatment of vector spaces in which feature dimensions, rather than only data points, participate in a structural semantic geometry. Given a corpus of embeddings with feature dimension $F$, one constructs a feature-space graph and its Laplacian $\Lf \in \R^{F \times F}$. The eigenvectors of $\Lf$ define smooth and rough directions in feature space, while associated bounded spectral scores quantify how individual items or signals inhabit this structure.  The more recent Energy Dispersion Networks formulation \cite{MoriondoAzizi2026} makes this perspective explicit by treating the feature columns of an embedding matrix as nodes in a semantic energy network. The resulting Laplacian spectrum and energy-dispersion statistics are not auxiliary diagnostics but a representation of how semantic structure is distributed over the feature graph.  In this paper we treat ArrowSpace as a frozen semantic prior: \[ \mathcal{G}_{\mathrm{Arrow}} = (\Lf,\Uq,\Lq,\ed),\] where $\Uq \in \R^{F \times q}$ contains the retained smooth modes, $\Lq = \mathrm{diag}(\nu_1,\ldots,\nu_q)$ are the corresponding eigenvalues, and $\ed \in [0,1]^F$ is an empirical energy-dispersion distribution over feature nodes.

\subsection{Latent diffusion}
\label{sec:latent_diffusion}

Latent diffusion models shift the forward and reverse diffusion processes from
pixel space into a learned compressed latent representation. An image
$x \in \R^{H \times W \times 3}$ is encoded into a latent tensor
$z \in \R^{h \times w \times c}$, denoised over a sequence of timesteps,
and decoded back to RGB. This substantially reduces computation compared to
pixel-space diffusion and is the basis of the latent-diffusion family.

Standard latent diffusion typically assumes that the VAE latent is sufficient
as the sole state on which the generative process must operate. Our claim is
that this is lossy (the presence of a missing feature-space information component is demonstrated in \cite{MoriondoAzizi2026}): a spatial latent is effective for local visual detail and a separate feature-space spectral state is needed to preserve global
semantic organisation under compression.

\subsection{Internal time and entropic motivation}
\label{sec:internal_time}

The original \ESDM~framing was motivated by the problem of time in closed
systems. In \cite{ConnesRovelli1994} it is argued that in a generally covariant theory, the
physical time-flow may arise from the thermodynamical state of the system
rather than from a universal external parameter. Barontini's recent cold-atom
experiment \cite{Barontini2026} provides an operational testbed for this idea: a closed system
partitioned into observed and unobserved sectors admits an entropic time
constructed from internal entropy exchange, and an effective Schr\"odinger-like
description parametrised by that internal time reproduces the measured
dynamics.

We retain this as a design principle: generative dynamics should be informed by
state-derived geometry and energy flow rather than only by an arbitrary external
counter. However, in the present architecture this principle is instantiated
through a frozen spectral prior and a heat-kernel chart, rather than through a
fully entropic recurrence at every step.

% ============================================================
\section{Design transition}
\label{sec:transition}
% ============================================================

Earlier \ESDM~drafts coupled five ingredients tightly: a dual-space semantic
operator (ArrowSpace), a wave recurrence over item states, a bright--dark sector entropy
clock, an entropic pump, and a topology decoder acting as an inverse spectral
map. This unified picture was conceptually rich but implementation-heavy.
Recent refactoring suggests a simpler separation of concerns.

\begin{remark}
The strongest reusable result of the earlier \ESDM~is not the full
second-order entropic dynamics but the geometric constraint that generation
should be restricted to a smooth ArrowSpace subspace and decoded relative to a
feature-space manifold according to an entropic diffusion scheduler.
\end{remark}

The current design therefore keeps:

\begin{itemize}[leftmargin=1.5em]
\item the frozen ArrowSpace prior;
\item the projection onto smooth spectral modes;
\item the discrete ("2.5-D") representation;
\item the heat-kernel interpretation of the spectral chart;
\item topology-adaptive decoding.
\end{itemize}

It defers for now as the next step for implementation:

\begin{itemize}[leftmargin=1.5em]
\item the full wave recurrence;
\item explicit Rayleigh-gradient restoring forces in the main forward pass;
\item the bright--dark KL clock as the primary training schedule;
\item the entropic pump as a learned forcing term.
\end{itemize}

This transition preserves the theoretical intent while producing a model that
is easier to train, analyse, and compare to ordinary latent diffusion systems. We present this model in the form of \ALDSC~\cite{Moriondo2026ALDSC_code}: a model that builds on top of \VDT~\cite{Moriondo2026VDT} and \ESDM~\cite{Moriondo2026ESDM_code} to deliver on concept of \emph{Spectral Latent Diffusion}. 

% ============================================================
\section{The 2.5-D latent}
\label{sec:25d}
% ============================================================

\subsection{Dual representation}
\label{sec:dual_representation}

Let $x \in \R^{H \times W \times 3}$ be an image. The encoder produces two
coupled outputs:
\[x \xrightarrow{E_\phi} (z, A),\]
where
\[z \in \R^{h \times w \times c}, \qquad A \in \R^{N \times F}, \qquad N = h \cdot w.\]

The tensor $z$ is the ordinary spatial latent used by the diffusion backbone.
The matrix $A$ is a feature field: each row corresponds to a spatial location
or patch, and each column corresponds to one feature dimension in the same
feature-space on which the ArrowSpace prior was built.

\subsection{Smooth projection}
\label{sec:smooth_projection}

The ArrowSpace prior defines a smooth subspace through the projection
\[\Pi_q = \Uq \Uq^\top.\]
Applying this to the feature field yields
\[A_{\parallel} = A \Pi_q, \qquad A_{\perp} = A(I - \Pi_q).\]
The parallel component lies in the smooth semantic subspace defined by the
lowest retained modes of $\Lf$, while the perpendicular component captures
high-frequency or off-manifold residual structure.

This gives the representation its "2.5-D" (or discrete-derived) character. The model preserves a full
spatial latent for local image geometry but constrains the semantic side of
the representation to a lower-dimensional graph manifold.

\subsection{Spectral chart}
\label{sec:spectral_chart}

From the projected feature field we extract a compact chart coordinate
\[s = \mathrm{Pool}(A_{\parallel}) \Uq \in \R^q,\]
together with per-mode band energies
\[e_k = \frac{1}{N}\|A u_k\|_2^2, \qquad k=1,\ldots,q,\]
where $u_k$ is the $k$-th retained eigenvector.

We normalise these energies by
\[\tilde e_k = \frac{e_k}{\sum_{j=1}^q e_j + \varepsilon}.\]
The projected energy-dispersion prior is
\[\ed_{k^{\mathrm{chart}}} = \sum_{f=1}^{F} \ed_f \, U_{f,k}^2.\]
The full spectral conditioning vector is then
\[c_{\mathrm{spec}} = \mathrm{MLP} \left( [\tilde e_1,\ldots,\tilde e_q,\; \ed_{1^{\mathrm{chart}}},\ldots,\ed_{q^{\mathrm{chart}}},\; \nu_1,\ldots,\nu_q] \right).\]


% ============================================================
\section{Entropic time from the ArrowSpace spectrum}
\label{sec:entropic_clock}
% ============================================================

The heat-kernel prior of Section~\ref{sec:heat_kernel} admits a natural
reading in terms of entropic time. We show that the ArrowSpace Laplacian
eigenvalues play the role of Barontini's barrier-controlled entropy exchange
rates, yielding a per-mode entropic clock that requires no additional learned
parameters and no auxiliary model. This connects the frozen spectral prior
directly to recent information-theoretic advances in diffusion scheduling \cite{Stancevic2025}.

% ------------------------------------------------------------
\subsection{Entropy rate of the heat-kernel chart}
\label{sec:entropy_rate}
% ------------------------------------------------------------

Recall the per-mode variance under the heat-kernel prior:
\[\sigma_k^2(t) = w_k \, e^{-2t\nu_k}, \qquad k = 1,\ldots,q,\]
where \(\nu_k\) is the \(k\)-th ArrowSpace Laplacian eigenvalue and \(w_k\)
is the projected energy-dispersion weight. The Shannon entropy of a
one-dimensional Gaussian with variance \(\sigma_k^2\) is
\[S_k(t) = \tfrac{1}{2}\log\!\bigl(2\pi e\, w_k\bigr) - t\,\nu_k.\]
The entropy rate is therefore
\[\boxed{\frac{dS_k}{dt} = -\nu_k.}\]
Each Laplacian eigenvalue \(\nu_k\) \emph{is} the entropy exchange rate for
mode~\(k\). Modes with small~\(\nu_k\) (smooth, globally coherent) exchange
entropy slowly; modes with large~\(\nu_k\) (rough, locally varying) exchange
entropy rapidly and reach equilibrium sooner.

% ------------------------------------------------------------
\subsection{Correspondence to Barontini's two-sector clock}
\label{sec:barontini_mapping}
% ------------------------------------------------------------

Barontini partitions a closed Bose--Einstein condensate into a bright
(observed) and dark (unobserved) sector, separated by a tunable potential
barrier of height~\(V\). Entropic time is constructed from the absolute
entropy exchange between sectors:
\[\tau = \int \left|\frac{dS_{\mathrm{bright}}}{d\phi}\right| \,|d\phi|,\]
where \(\phi\) is a clock field. The barrier height controls the rate of
entropy production and thus the speed of the emergent internal time.

\ESDM, as it is defined in 

Table~\ref{tab:barontini} maps this construction onto the \ALDSC~spectral
chart.

\begin{table}[ht]
\centering
\small
\caption{Correspondence between Barontini's BEC experiment and the ALD-SC
spectral chart.}
\label{tab:barontini}
\begin{tabular}{ll}
\toprule
\textbf{Barontini experiment} & \textbf{ALD-SC spectral chart} \\
\midrule
Bright sector (observed) & Active modes: low \(\nu_k\), high variance \\
Dark sector (unobserved) & Inactive modes: high \(\nu_k\), collapsed variance \\
Barrier height \(V\) & Laplacian eigenvalue \(\nu_k\) (per-mode barrier) \\
Entropy exchange rate & \(|dS_k/dt| = \nu_k\) \\
Entropic time \(\tau\) & Per-mode \(\tau_k(t) = \nu_k \cdot t\) \\
Big bang / big crunch & Mode activation / collapse during denoising \\
Heat death & All modes at equilibrium variance \\
Barrier-height sweep & Varying \(q\) (number of retained modes) \\
\bottomrule
\end{tabular}
\end{table}

The barrier-height sweep that Barontini uses to control the speed of entropic
time maps to the choice of~\(q\). Fewer retained modes means higher average
barrier and slower, more conservative diffusion; more modes means lower
average barrier and faster, more permissive exchange.

% ------------------------------------------------------------
\subsection{Per-mode entropic schedule}
\label{sec:per_mode_schedule}
% ------------------------------------------------------------

Instead of a single global schedule \(\bar\alpha(t)\), each mode receives its
own entropic schedule driven by its internal clock:
\[\tau_k(t) = \nu_k \cdot t, \qquad \bar\alpha_k(\tau_k) = \cos^2\!\left(\frac{\pi}{2}\cdot\frac{\tau_k}{\tau_{k,\max}}\right),\]
where \(\tau_{k,\max} = \nu_k \cdot T\) for a reference horizon~\(T\).
The external time at which mode~\(k\) reaches heat death is
\[t_k^{\mathrm{death}} = \frac{\tau_{k,\max}}{\nu_k} = T.\]
While all modes formally reach equilibrium at the same external time~\(T\),
they traverse their entropic trajectories at different rates. High-\(\nu_k\)
modes accumulate entropic time faster, so they pass through the critical
transition regime (where \(\bar\alpha_k\) crosses~\(1/2\)) earlier in the
denoising process. Low-\(\nu_k\) modes reach this transition later.

\begin{remark}
This yields a topology-adaptive ordering: global semantic structure (low
\(\nu_k\)) persists throughout the denoising trajectory, while local detail
(high \(\nu_k\)) is resolved or lost according to the graph geometry --- not
by design choice but by the spectral structure of the feature-space graph.
\end{remark}

% ------------------------------------------------------------
\subsection{Connection to information-theoretic schedulers}
\label{sec:stancevic}
% ------------------------------------------------------------

Stancevic et al.~\cite{Stancevic2025} prove that the conditional entropy
\(\mathbf{H}[\mathbf{x}_0 \mid \mathbf{x}_t]\) is a schedule-invariant time
reparameterisation for diffusion models, and provide a tractable estimator
from the training loss via the I-MMSE identity. The entropy rate decomposes
over modes when the chart covariance is diagonal:
\[\dot{\mathbf{H}}[s_0 \mid s_t] = \sum_{k=1}^{q} \nu_k \cdot \mathrm{mmse}_k(t),\]
where \(\mathrm{mmse}_k(t)\) is the per-mode minimum mean-square error of
recovering \(s_0\) from~\(s_t\). The total entropic time is
\[\tau(t) = \int_0^t \sum_{k=1}^{q} \nu_k \cdot \mathrm{mmse}_k(t')\, dt'.\]
Stancevic et al.~\cite{Stancevic2025} show that this quantity can be estimated from the unweighted
denoising loss without auxiliary models. The \ALDSC~contribution is that the
ArrowSpace spectrum provides the decomposition basis a priori: the
\(\nu_k\) are frozen graph eigenvalues conditioned on the training set (via the ArrowSpace learned structure).

Table~\ref{tab:scheduler_comparison} contrasts \ALDSC~with prior
information-theoretic and spectral scheduling approaches.

\begin{table}[ht]
\centering
\small
\caption{Comparison of noise-scheduling approaches.}
\label{tab:scheduler_comparison}
\begin{tabular}{llll}
\toprule
\textbf{Feature} & \textbf{Stancevic} & \textbf{Spectrally-Guided} & \textbf{ALD-SC} \\
 & \textbf{et al.} & \textbf{(Esteves)} & \textbf{(this work)} \\
\midrule
Time variable & Cond.\ entropy & RAPSD profile & Per-mode \(\tau_k\) \\
Spectral basis & Fourier (pixel) & Fourier (pixel) & ArrowSpace (semantic) \\
Per-mode schedule & Spectral variant & No & Yes, by construction \\
Geometry source & Image frequency & Image spectrum & Feature-space graph \\
Online adaptation & No (post-train) & No & Yes (MMSE online) \\
Barrier analogue & None & None & \(q\) (retained modes) \\
\bottomrule
\end{tabular}
\end{table}

% ------------------------------------------------------------
\subsection{Training: entropic noise allocation}
\label{sec:training_clock}
% ------------------------------------------------------------

Following the information-guided noise allocation principle, the per-mode
information rate is estimated online as
\[\hat\rho_k(t) \propto \nu_k \cdot \hat{m}_k(t),\]
where \(\hat{m}_k(t)\) is the running average of the per-mode denoising loss.
The training noise distribution is adapted so that
\[\pi(t) \propto \frac{\hat\rho(t)}{w(t)},\]
concentrating training effort at noise levels where spectral modes are
resolving information fastest. This operationalises the Barontini principle
that time advances fastest where entropy is flowing most rapidly.

A warm-up period with a fixed cosine schedule is recommended during early
training, when per-mode losses are noisy. The entropic clock is activated
once the running loss estimates stabilise.

% ------------------------------------------------------------
\subsection{Inference: entropic step placement}
\label{sec:inference_clock}
% ------------------------------------------------------------

At sampling time, steps are placed equidistantly in the rescaled entropic time
\[\tau_{\mathrm{rescaled}}(t) = \int_0^t \sigma(\tau)\, \dot{\mathbf{H}}[s_0 \mid s_\tau]\, d\tau,\]
rather than in external time~\(t\). This ensures each step contributes equal
information, with more steps allocated to the critical windows where modes
are actively resolving.

\begin{remark}
Stancevic et al.~report that entropic time schedules perform best with
first-order solvers (DDIM, flow matching). Second-order solvers may require
adjustment because the lookahead step changes the effective entropy rate.
\end{remark}

% ------------------------------------------------------------
\subsection{Heat death as a stopping criterion}
\label{sec:heat_death_clock}
% ------------------------------------------------------------

Barontini observes that entropic time stops advancing when no entropy is
exchanged between sectors. In \ALDSC, this corresponds to all modes reaching
equilibrium variance:
\[\mathrm{heat\ death} \iff \sum_{k=1}^{q} \nu_k \cdot \mathrm{mmse}_k(t) < \varepsilon.\]
This provides an intrinsic stopping criterion for sampling: the denoising
process terminates when no further information is being resolved by any
spectral mode, rather than at a fixed step count. The number of effective
sampling steps is therefore data-dependent and spectrally determined.

% ------------------------------------------------------------
\subsection{Summary}
\label{sec:clock_summary}
% ------------------------------------------------------------

The ArrowSpace Laplacian spectrum provides a per-mode entropic clock that is
derived from the same frozen geometric prior that defines the semantic
manifold. No additional learned parameters are required: the eigenvalues
\(\nu_k\) are computed once from the corpus feature graph and remain fixed
throughout training and inference. The schedule is not externally prescribed
but emerges from the interaction between the graph geometry and the current
denoising state, connecting the practical architecture to the stronger
thermal-time programme without entangling the implementation.



% ============================================================
\section{Model architecture}
\label{sec:model}
% ============================================================

\subsection{Overview}
\label{sec:overview}

The revised \ESDM~(\ALDSC)~consists of four components:

\begin{enumerate}[leftmargin=1.5em]
\item a frozen ArrowSpace prior $(\Lf,\Uq,\Lq,\ed)$;
\item a spectral VAE encoder producing a spatial latent $z$ and feature field $A$;
\item a latent diffusion denoiser operating only on $z$;
\item a topology-adaptive decoder conditioned on the spectral chart.
\end{enumerate}

The full pathway is
\[x \;\xrightarrow{E_\phi}\; (z, A) \;\xrightarrow{\text{chart}} (z, c_{\mathrm{spec}}) \;\xrightarrow{\text{diffusion}} \hat z \;\xrightarrow{D_\psi(\cdot \mid c_{\mathrm{spec}})} \hat x.\]

\subsection{Encoder}
\label{sec:encoder}

The encoder is a shared visual trunk with two heads. The first head produces
the spatial latent
\[z \sim q_\phi(z \mid x),\]
typically via a mean and log-variance field as in a standard VAE. The second
head produces the feature field
\[A = E_A(x) \in \R^{N \times F}.\]
The feature field is not itself diffused; it is projected onto the ArrowSpace
subspace and summarised into spectral tokens.

\subsection{Diffusion backbone: the spectral chart}
\label{sec:backbone}

We apply diffusion only to the spatial latent:
\[q(z_t \mid z_0) = \mathcal{N} \left( \sqrt{\bar\alpha_t} z_0,\; (1-\bar\alpha_t)I \right).\]
A denoiser predicts velocity or noise:
\[v_\theta(z_t, t, c_{\mathrm{text}}, c_{\mathrm{spec}}).\]
The spectral chart enters the denoiser through conditioning channels such as
cross-attention or adaptive layer normalisation. In this way, the diffusion
backbone remains computationally standard, while semantic topology enters as a
side-channel rather than as a second full state to denoise.

\subsection{Topology-adaptive decoder}
\label{sec:decoder}

The decoder reconstructs pixels from the denoised spatial latent under
spectral control:
\[\hat x = D_\psi(\hat z; c_{\mathrm{spec}}).\]
At each resolution stage the decoder forms feature activations, projects them
onto the chart basis $\Uq$, gates these modes according to
$c_{\mathrm{spec}}$, and lifts them back to feature space. Abstractly, a
decoder block has the form
\[H^{(\ell+1)} = \mathrm{UpBlock} \left( H^{(\ell)} + (\hat H^{(\ell)} \odot g^{(\ell)})\Uq^\top \right),\]
where $\hat H^{(\ell)} = H^{(\ell)}\Uq$ and
$g^{(\ell)} = \sigma(W_\ell c_{\mathrm{spec}} + b_\ell)$.

This construction makes the decoder topology-adaptive: \emph{reconstruction is not
performed relative to an isotropic ambient prior alone but relative to the
smooth directions described by the ArrowSpace graph}.

% ============================================================
\section{Heat-kernel prior}
\label{sec:heat_kernel}
% ============================================================

The spectral chart admits a natural diffusion-geometric prior. Let
$s = (s_1,\ldots,s_q)^\top$ denote the chart coordinates. We define
\[p_t(s) = \mathcal{N} \left( 0,\; \mathrm{diag}\bigl[w_k e^{-2t\nu_k}\bigr]_{k=1}^q \right),\]
where the weights $w_k$ are derived from the projected energy-dispersion prior:
\[w_k = \frac{\exp(-\rho \ed_{k^{\mathrm{chart}}})} {\sum_j \exp(-\rho \ed_{j^{\mathrm{chart}}})}.\]

This prior realises the heat-kernel view of the ArrowSpace manifold. Low
eigenvalue modes persist under diffusion and therefore encode broad semantic
structure; high eigenvalue modes decay more rapidly and correspond to rough or
unstable directions. The energy-dispersion weights bias the chart toward the
regions of spectral space actually occupied by corpus data.

\begin{definition}[2.5-D manifold prior]
The pair $(z,s)$ defines a 2.5-D manifold prior if $z$ carries local spatial
degrees of freedom and $s$ is drawn from a heat-kernel distribution on the
ArrowSpace spectral chart.
\end{definition}

% ============================================================
\section{Objective}
\label{sec:objective}
% ============================================================

The revised training objective is
\[\mathcal{L} = \mathcal{L}_{\mathrm{diff}} + \lambda_{\mathrm{rec}}\mathcal{L}_{\mathrm{rec}} + \lambda_{\mathrm{chart}}\mathcal{L}_{\mathrm{chart}} + \lambda_{\mathrm{smooth}}\mathcal{L}_{\mathrm{smooth}} + \lambda_{\mathrm{kl}}\mathcal{L}_{\mathrm{kl}}.\]

\subsection{Latent diffusion loss}
\label{sec:loss_diff}

Using velocity prediction,
\[\mathcal{L}_{\mathrm{diff}} = \mathbb{E}_{z_0,\epsilon,t} \left[ \left\| v - v_\theta(z_t,t,c_{\mathrm{text}},c_{\mathrm{spec}}) \right\|_2^2 \right].\]

\subsection{Reconstruction and VAE regularisation}
\label{sec:loss_rec}

The image reconstruction term is
\[\mathcal{L}_{\mathrm{rec}} = \|x-\hat x\|_1 + \lambda_{\mathrm{perc}}\mathcal{L}_{\mathrm{perc}}(x,\hat x),\]
and the VAE regulariser is the usual latent KL:
\[\mathcal{L}_{\mathrm{kl}} = D_{\mathrm{KL}}\bigl(q_\phi(z \mid x)\,\|\,p(z)\bigr).\]

\subsection{Spectral consistency}
\label{sec:loss_chart}

To preserve global semantic organisation, the spectral-band allocation of the
decoded image should match that of the source image. Re-encoding $\hat x$
produces $\hat A$, and we define
\[\mathcal{L}_{\mathrm{chart}} = \left\| \tilde e(x) - \tilde e(\hat x) \right\|_2^2.\]

\subsection{Off-manifold penalty}
\label{sec:loss_smooth}

We penalise feature energy outside the smooth ArrowSpace subspace:
\[\mathcal{L}_{\mathrm{smooth}} = \frac{ \|A(I-\Uq\Uq^\top)\|_F^2 }{ \|A\|_F^2 + \varepsilon }.\]

This is the cleanest operational descendant of the earlier \ESDM claim that
generation should be restricted to semantically coherent modes.

% ============================================================
\section{Experimental programme}
\label{sec:experiments}
% ============================================================

We propose a staged empirical evaluation.

\subsection{Stage 1: spectral VAE}
\label{sec:stage1}

Train the spectral VAE without diffusion and compare it to a conventional VAE
at equal latent capacity. The main questions are:

\begin{itemize}[leftmargin=1.5em]
\item Does spectral conditioning improve reconstruction fidelity?
\item Does it reduce off-manifold energy?
\item Does it preserve more stable global semantic organisation under
interpolation and compression?
\end{itemize}

\subsection{Stage 2: full latent diffusion}
\label{sec:stage2}

Train the latent denoiser on the spatial latent while conditioning on the
spectral chart. Compare against a baseline latent diffusion model without
ArrowSpace conditioning.

\subsection{Metrics}
\label{sec:metrics}

Beyond standard image-quality metrics, evaluation should include spectral
diagnostics:

\begin{itemize}[leftmargin=1.5em]
\item chart-energy error;
\item off-manifold energy ratio;
\item chart interpolation smoothness;
\item reconstruction agreement under re-encoding;
\item global semantic coherence under low latent capacity.
\end{itemize}

The central claim is not merely improved image quality but improved
semantic coherence under compression.

% ============================================================
\section{Connection to internal time}
\label{sec:internal_time_connection}
% ============================================================

Although the current model does not use the full entropic clock as its primary
training schedule, the original motivation remains important. The conceptual
bridge is that geometry and dynamics should be derived from internal state
rather than imposed entirely from outside. In the original \ESDM~this took the
form of an entropy-derived temporal variable. In the present architecture it
takes the form of a feature-space geometric prior that governs which semantic
directions are valid and how their energy is distributed.

This suggests a future synthesis: the spectral chart could itself support a
state-derived adaptive schedule, for example by coupling noise level to
off-manifold energy, chart entropy, or spectral-band redistribution. Such a
development would reconnect the practical architecture to the stronger thermal
time programme without entangling the initial implementation.

% ============================================================
\section{Conclusion}
\label{sec:conclusion}
% ============================================================

We presented a revised Entropic Semantic Diffusion Model in which the main
innovation is not a fully endogenous clock but a topology-adaptive latent
diffusion architecture grounded in ArrowSpace feature geometry. The model
combines a standard spatial latent with a frozen spectral prior and a decoder
conditioned on a compact ArrowSpace chart. This realises the central 2.5-D
claim: local image detail lives in a spatial latent plane, while global
semantic coherence is governed by a lower-dimensional graph manifold.

The original entropic formulation remains an important theoretical horizon.
However, the present architecture offers a cleaner and more tractable route to
testing the key scientific idea: that generative models should reconstruct and
sample not only in compressed latent coordinates but on a semantic manifold
defined by the graph Laplacian structure of the feature space itself.

% ============================================================
\newpage
\begin{thebibliography}{99}

\bibitem{Moriondo2026VDT}
L.~ Moriondo,
\newblock {Vibrational Deduction Transformer (VDT) Leveraging Spectral methods, Theory of Sound and Associative Memory}.
\newblock Zenodo.
\newblock \url{https://doi.org/10.5281/zenodo.20816835}.


\bibitem{Barontini2026}
G.~Barontini,
\newblock {Testing the problem of time with cold atoms},
\newblock {\em Physical Review Research}, 8:L022047, 2026.
\newblock doi: \href{https://doi.org/10.1103/1h9j-df4k}{10.1103/1h9j-df4k}.

\bibitem{ConnesRovelli1994}
A.~Connes and C.~Rovelli,
\newblock {Von Neumann Algebra Automorphisms and Time-Thermodynamics Relation in General Covariant Quantum Theories},
\newblock {\em Classical and Quantum Gravity}, 11:2899--2918, 1994.
\newblock \href{https://arxiv.org/abs/gr-qc/9406019}{arXiv:gr-qc/9406019}.

\bibitem{Moriondo2025JOSS}
L.~Moriondo,
\newblock {ArrowSpace: introducing Spectral Indexing for vector search},
\newblock {\em Journal of Open Source Software}, 10(113):9002, 2025.
\newblock doi: \href{https://doi.org/10.21105/joss.09002}{10.21105/joss.09002}.

\bibitem{MoriondoAzizi2026}
L.~Moriondo and I.~Azizi,
\newblock {From Embedding Geometry to Spectral Search: Energy Dispersion Networks for Vector Retrieval},
\newblock arXiv:2606.21535 [cs.IR], 2026.
\newblock \href{https://arxiv.org/abs/2606.21535}{https://arxiv.org/abs/2606.21535}.

\bibitem{Rombach2022}
R.~Rombach, A.~Blattmann, D.~Lorenz, P.~Esser, and B.~Ommer,
\newblock {High-Resolution Image Synthesis with Latent Diffusion Models},
\newblock in {\em CVPR}, 2022.
\newblock \href{https://arxiv.org/abs/2112.10752}{arXiv:2112.10752}.

\bibitem{NicholDhariwal2021}
A.~Nichol and P.~Dhariwal,
\newblock {Improved Denoising Diffusion Probabilistic Models},
\newblock in {\em ICML}, 2021.
\newblock \href{https://arxiv.org/abs/2102.09672}{arXiv:2102.09672}.

\bibitem{Moriondo2026ESDM_code}
L.~Moriondo,
\newblock {\textsc{\Large\faGithub} entropic-semantic-diffusion},
\newblock \href{https://github.com/tuned-org-uk/entropic-semantic-diffusion}{Github}

\bibitem{Moriondo2026ALDSC_code}
L.~Moriondo,
\newblock {\textsc{\Large\faGithub} arrowspace-latent-diffusion},
\newblock \href{https://github.com/tuned-org-uk/arrowspace-latent-diffusion}{Github}

\bibitem{Stancevic2025}
D.~Stancevic, F.~Handke, and L.~Ambrogioni,
\newblock {Entropic Time Schedulers for Generative Diffusion Models},
\newblock arXiv:2504.13612 [cs.LG], 2025.
\newblock \href{https://arxiv.org/abs/2504.13612}{arXiv:2504.13612}.

\bibitem{Esteves2026}
C.~Esteves and A.~Makadia,
\newblock {Spectrally-Guided Diffusion Noise Schedules},
\newblock in {\em ICML}, 2026.
\newblock \href{https://arxiv.org/abs/2603.19222}{arXiv:2603.19222}.

\end{thebibliography}

\end{document}
